\subsection{Unit-root eighth moments for terminal pairs \((3,b)\)}
\label{proof:Q5-C089}
Keep the finite model and notation of Section~\ref{proof:Q5-C050}, and
write \(f_n=A^n\one\), \(h_n=Tf_{n-1}\),
\(\bar f_n=T^n\one\). Choose the terminal direction from the integer
lengths before choosing the host:
\begin{equation}\label{eq:unit-ordinary:three-directions}
\begin{array}{c|c}
 b\ge4, L=1&f_3h_b\\
 b=4, L=2&h_3f_4\\
 b\ge5, L=2&f_3h_b\\
 b\ge4, L\ge3&h_3f_b.
\end{array}
\end{equation}
An \(f_n\) arm points outward throughout. In an \(h_n\) arm only
its first edge points toward the terminal. The connector points from
the terminal to the receiving root, whose original unit arm supplies
\(p\). For actual equal-pair factors \(0\le E_t,E_r\le1\), set
\begin{equation}\tag{Q89.0.2}\label{eq:Q89:0.2}
 H_D=E_rpT^L(E_tF_{b,L}).
\end{equation}
We prove
\begin{equation}\tag{Q89.0.3}\label{eq:Q89:0.3}
 \E H_D^8\le1-d/500\qquad(b\ge4, L\ge1).
\end{equation}
The \(L=1\) row is Section~\ref{proof:Q5-C053-unit}. Positivity
allows us to set \(E_t=E_r=1\) until the end. At \(d=0\), all
tree messages are one by leaf integration.

\paragraph{Finite square propagation at the original unit root.}
Put \(K'=-K\), \(g'=K'\one\), \(G_n(z)=\sum_{j=0}^{n-1}z^j\),
\(G_0=0\). The finite convolution \eqref{eq:Q188:conv} gives
\[
 \tau\le\beta_j\le d\ (j\ge3),\qquad
 \|G_n(K')\|\le1\ (n=0,1\text{ or }n\ge3),\quad
 \|G_2(K')\|\le\sqrt{1+\theta}<6/5.
\]
Suppose \(0\le F\le(m+R/2)^2\), where \(m>0\), and put
\[
\begin{gathered}
 k=m^{-1},\quad D_j=\E(K')^jR,\quad V=\E R^2,
 \quad a_1=-D_1,\quad\rho=-D_2,\\
 \delta=kD_0+k^2V/2,\quad a=ka_1,\quad\varepsilon=\delta-a.
\end{gathered}
\]
Expanding the square, using \(T(R^2)\le2V\), and applying finite
Duhamel gives
\begin{equation}\label{eq:unit-ordinary:three-duhamel}
 \frac{T^LF}{m^2}\le
 \bar f_L+k(K')^LR+\delta\bar f_{L-1}-a\bar f_{L-2}
       +k\sum_{j=2}^{L-1}D_j\bar f_{L-1-j}\quad(L\ge2).
\end{equation}
For \(R=(A(x)-B(x)z)g'\), \(z=i\sqrt x\), use the same
spectral measure \(\nu\) of \(g'\), supported on \([0,\theta]\)
and of mass \(d\). Thus
\begin{equation}\label{eq:unit-ordinary:three-moments}
 \begin{gathered}
 V=\int(A^2+xB^2)\,d\nu,\quad D_0=\int B\,d\nu,
 \quad a_1=\int A\,d\nu,\quad\rho=\int xB\,d\nu,\\
 D_{4j+1}=-\int x^{2j}A\,d\nu,\quad
 D_{4j+2}=-\int x^{2j+1}B\,d\nu,\\
 D_{4j+3}=\int x^{2j+1}A\,d\nu,\quad
 D_{4j+4}=\int x^{2j+2}B\,d\nu.
 \end{gathered}
\end{equation}
Let \(S'_L=\sum_{j=0}^L(K')^j\one\), and define
\[
 \begin{gathered}
 P_n=G_n+kz^n(A-Bz),\quad
 Q_n=P_n+\varepsilon G_{n-2}+\delta z^{n-2}\quad(n\ge2),\\
 Z_L^*=S'_L+\delta S'_{L-1}-aS'_{L-2}+k(K')^LR
       =1+\varepsilon+Q_L(K')g'.
 \end{gathered}
\]
The exact difference between \(Z_L^*\) and the right side of
\eqref{eq:unit-ordinary:three-duhamel} is
\begin{equation}\label{eq:unit-ordinary:three-losses}
 \begin{gathered}
 \sum_{j=3}^L\gamma_j\bar f_{L-j},\qquad
 \gamma_3=d+k\rho,\quad\gamma_4=d+\delta d-kD_3,\\
 \gamma_j=\beta_j+\delta\beta_{j-1}-a\beta_{j-2}-kD_{j-1}
 \qquad(j\ge5).
 \end{gathered}
\end{equation}
Also \(Q_L=1+\varepsilon+zQ_{L-1}\) for \(L\ge3\), and
\(\E Z_L^*=1+\varepsilon-\gamma_{L+1}\) for \(L\ge4\).
These are finite identities, including the constant coordinate.

Whenever \(N=K'W\), \(0\le H\le(c+N/2)^2\), and the square
base is nonnegative, the signed estimate of
Section~\ref{proof:Q5-C081} and \eqref{eq:oc:ordinary-square} give
\begin{equation}\label{eq:unit-ordinary:three-signed}
 \E H^8\le\exp\{16(c-1)+8\E N+17\|W\|^2\}.
\end{equation}
For a normalized square with prefactor \(m^{16}\) this exponent
becomes \(16\log m+8\eta+8\E N+17\|W\|^2\), when its center
before rescaling is \(1+\eta/2\).

\paragraph{Two terminal packets, with all finite arm lengths.}
For this paragraph assume \(0<d\le1/5\). At \(b=4\), the exact
identity \(h_3+f_4=2-3d+t-dg\) gives
\begin{equation}\label{eq:unit-ordinary:three-four-sum}
 \begin{gathered}
 h_3+f_4=2m+R,\qquad m=1-3d/2,\\
 R=t-dg,\qquad A=d,\quad B=x.
 \end{gathered}
\end{equation}
In particular
\begin{equation}\label{eq:unit-ordinary:three-four-data}
 \begin{gathered}
 D_0=\delta_2,\quad a_1=d^2,\quad\rho=\delta_3,
 \quad V=d^3+\delta_4,\\
 U:=\|KR\|^2=d^2\delta_2+\delta_5,\qquad
 W:=\|K^2R\|^2=d^2\delta_3+\delta_6.
 \end{gathered}
\end{equation}
Moreover \(0\le\varepsilon\le\delta\le7d/10\), as will also
follow from the scalar certificates below.

For every \(b\ge5\), retain three terms of each positive convolution.
At \(b=5\) take \(c_1=d\), \(c_2=d/4+\delta_2\), \(c_3=0\).
For \(b\ge6\), take
\begin{equation}\label{eq:unit-ordinary:three-retained}
 \lambda=\frac{9d/4-\beta_b-\beta_{b-1}}
                 {\beta_{b-2}+(1-d)\beta_{b-3}},\quad
 c_1=\beta_{b-1},\quad c_2=\lambda\beta_{b-2},\quad
 c_3=\lambda\beta_{b-3}.
\end{equation}
The denominator is at least \((2-d)\tau>0\), the numerator is
nonnegative, and the available constant loss is at least
\((4-d)\tau\ge361d/160>9d/4\); hence \(0\le\lambda\le1\).
For \(b=5\), \(c_2\le21d/32\le\beta_3\), without using
the zero coefficient \(\beta_2\).
Put \(C=c_1+c_2+c_3\), \(D=c_2+c_3\), \(E=c_3\), and
\[
 R_{\rm old}=\sum_{j=3}^bK^j\one-Cg-Dr-Es,\quad
 R=R_{\rm old}-s,\qquad m=1-13d/8.
\]
The two positive convolutions yield the two actual arm sums
\begin{equation}\label{eq:unit-ordinary:three-arm-sums}
 \begin{gathered}
 h_2+f_b\le2-9d/4+R_{\rm old},\quad
 f_2+h_b\le2-9d/4-R_{\rm old},\\
 h_3+f_b\le2m+R,\qquad f_3+h_b\le2m-R.
 \end{gathered}
\end{equation}
Both signs therefore give positive square envelopes. For the second
sign retain the actual \(q,h_2,h_3\) factors after applying \(T\).
Their constant loss in either direction is
\(\beta_b+c_1+c_2+(1-d)c_3=9d/4\).

Here is the complete finite coordinate description. At \(b=5\),
\(A=C-x^2\), \(B=D+x\). At \(b\ge6\), put \(N=b-2\ge4\)
and \(G_N=a_f+zb_f\). Then
\begin{equation}\label{eq:unit-ordinary:three-coordinates}
 A=C+x(a_f-E-1),\quad B=D+xb_f,\quad
 a_f,b_f\in[1-x,1-x+x^2],\quad a_f\ge(1-x)b_f.
\end{equation}
These finite geometric bounds include \(N=4,5\).
The retained constants satisfy
\[
 \begin{gathered}
 C=9d/4-\beta_b+dE,\quad
 D=9d/4-\beta_b-\beta_{b-1}+dE,\quad E\le d,\\
 \tau\le c_1\le d,\quad C\le d+D,\quad
 d/4\le D\le85d/64,\quad C\le C_0:=(53/32+d)d.
 \end{gathered}
\]
Consequently \(-x^2\le A\le C\),
\(D+(1-x)x\le B\le D+x\), and
\(A+x-(1-x)B\ge c_1\ge\tau\).
Using \(D\le d/4+2\delta_2+dE\) in the upper bound for \(B\)
gives
\[
 B\le x-x^2+x^3+d/4+2\delta_2+dE
   \le z_0:=\theta+13d/16.
\]
Indeed the difference after setting \(x=\theta\),
\(\delta_2\le\theta d\), \(E\le d\) is at most
\(-3/32+d/4+d^2\le-3/800\); the \(b=5\) bound is smaller.
Thus, with
\begin{equation}\label{eq:unit-ordinary:three-joint-data}
 v_0=C_0^2+\theta^4+\theta z_0^2,\quad
 s_0=kz_0+\tfrac12k^2v_0,\quad
 e_0=k[\theta(1+z_0)-\kappa d]+\tfrac12k^2v_0,
\end{equation}
we have \(V\le v_0d\), \(\delta\le s_0d\), and
\(\varepsilon\le e_0d\). The last bound uses the same spectral
measure and \(D_0-a_1\le[\theta(1+z_0)-\kappa d]d\).
For the opposite bound, \(a_f-b_f\) is zero or an alternating
end term; its first positive possibility is \(x^2\). Hence
\[
 a_1-D_0\le d^2-\delta_2+\delta_4\le\theta^3d,
 \qquad\varepsilon\ge-k\theta^3d>-d/10.
\]
At \(b=5\), the left side is exactly \(d^2-\delta_2-\delta_3\).

One additional pointwise estimate is needed for the exchanged
two-step direction:
\begin{equation}\label{eq:unit-ordinary:three-relative}
 A^2+xB^2\le(5/4)B<2mB.
\end{equation}
Indeed its ratio to \(B\) is at most
\((C^2+x^4)/(D+\kappa x)+x(D+x)\), a convex function of \(x\).
At \(x=0\) it is at most \(d(1+D/d)^2/(D/d)\le5/4\).
At \(x=\theta\), replace \(C\) by \(d+D\); the resulting function
is convex in \(D\), and then in \(d\) on each endpoint
\(D/d\in\{1/4,85/64\}\). The corner bounds are
\(169/608\), \(8/15\), and \(13/16\), all below \(5/4\).
Also \(2m\ge27/20>5/4\).

\paragraph{A single continuous payment for each short connector.}
For the canonical square, direct AM--GM gives the following two
complete vector formulas:
\begin{equation}\label{eq:unit-ordinary:three-short-vectors}
\begin{gathered}
 V_2=\delta\one+g'+kK'R,\qquad
 \E(K'V_2)=-d-k\rho,\\
 \|V_2\|^2=d+2k\rho+k^2U+\delta^2-2\delta a,\\
 V_3=\varepsilon\one+(1+\delta)g'+K'g'+k(K')^2R,\\
 \E(K'V_3)=-(1+\delta)d+kD_3,\\
 \begin{aligned}
 \|V_3\|^2={}&(1+\delta)^2d+\delta_2-2k(1+\delta)D_3\\
 &+2kD_4+k^2W+\varepsilon^2-2\varepsilon(d+k\rho).
 \end{aligned}
\end{gathered}
\end{equation}
For \(L=3\) the one retained path loss is \(\gamma_3=d+k\rho\).
The square bases are halves of sums of nonnegative factors. Applying
\eqref{eq:unit-ordinary:three-signed} yields
\begin{equation}\label{eq:unit-ordinary:three-short-payments}
\begin{aligned}
 \log\E(pT^2F)^8\le{}&16\log m+9d+8\varepsilon+26k\rho
                  +17k^2U+17(\delta^2-2\delta a),\\
 \log\E(pT^3F)^8\le{}&16\log m+(253/32)d+26\delta d+17\delta^2d\\
 &-(26+34\delta)kD_3+34kD_4-8k\rho+17k^2W+\mathcal E,\\
 \mathcal E={}&8\varepsilon+17\varepsilon^2
                         -34\varepsilon(d+k\rho).
\end{aligned}
\end{equation}
In particular the full constant-vector cross terms remain in both rows.

Here are explicit rational functions which pay these bounds on the
single interval \(0\le d\le1/5\). Write
\(\mathcal L_q=-16\sum_{j=1}^4(qd)^j/j\ge16\log(1-qd)\).
For \(b=4\), let \(m_4=1-3d/2\), \(k_4=1/m_4\), and
\(s_4=k_4\theta+\frac12k_4^2(d^2+\theta^3)\). Define
\[
\begin{split}
 F_{42}={}&\mathcal L_{3/2}+9d+8s_4d-8k_4d^2\\
 &+26k_4\theta^2d+17k_4^2(d^2\theta+\theta^4)d
 +17s_4^2d^2-34s_4k_4d^3,\\
 F_{43}={}&\mathcal L_{3/2}+(253/32)d+(91/5)d^2+(833/100)d^3\\
 &+34k_4\theta^3d+17k_4^2(d^2\theta^2+\theta^5)d
 +(28/5)d-(1547/100)d^2.
\end{split}
\]
The two bounds are \(F_{42}\le-d\), \(F_{43}\le-6d\).
For the first, the dependence on \(\delta\) is increasing and
\(\delta\le s_4d\); for the second use
\(0\le\varepsilon\le\delta\le7d/10\), \(D_3\ge0\),
and the data \eqref{eq:unit-ordinary:three-four-data}.

For \(b\ge5\), use \(m=1-13d/8\), \(k=1/m\), and the joint
functions \eqref{eq:unit-ordinary:three-joint-data}. Put
\[
 \begin{gathered}
 l=k\theta^3,\quad j_0=1+k\theta z_0,\quad
 E_+=8e_0d+(17e_0^2-34e_0)d^2,\quad
 E_-=-8ld+(17l^2+34j_0l)d^2,\\
 \begin{aligned}
 F_{\ge5,3}={}&\mathcal L_{13/8}+(253/32)d+26s_0d^2+17s_0^2d^3\\
 &+(26+34s_0d)k\theta^3d+34k\theta^2z_0d
                       +17k^2\theta^2v_0d+E_+.
 \end{aligned}
 \end{gathered}
\]
Convexity in \(\varepsilon\), retaining its sign, gives
\(\mathcal E\le\max(0,E_+,E_-)=E_+\).
Indeed \(-ld\le\varepsilon\le e_0d\),
\(d\le d+k\rho\le j_0d\), and the certificates below show
\(E_+\ge0,E_-\). The remaining substitutions are
\(D_3\ge-\theta^3d\), \(D_4\le\theta^2z_0d\),
\(W\le\theta^2v_0d\). They prove the whole-domain bound
\(F_{\ge5,3}\le-d/50\).

For the actual exchanged \(L=2\) packet replace \(R\) by \(-R\).
Equation~\eqref{eq:unit-ordinary:three-relative} gives
\(\delta_-=-kD_0+k^2V/2\le0\) and
\(\|g'-kK'R\|^2=d-2k\rho+k^2U\le d\), since
\(U\le(5/4)\rho\), \((5/4)k<2\).
The actual mean is \(-d+k\rho\), and
\[
 |\delta_-|\le6d/7,\quad |a_-|\le4d/7,
 \quad\varepsilon_-\le d/2,\quad k\rho\le7d/20,
 \quad\delta_-^2-2\delta_-a_-\le12d^2/7.
\]
Thus its exponent is at most
\([-26+9+4+14/5+204/35]d=-153d/35<-4d\).

The continuous certificates just used have only 78 rational
Bernstein coefficients in total. Each polynomial below has all its
coefficients greater than \(1/200\) on \([0,1/5]\).
The positive factors in the second column clear all denominators.
At \(d=0\), quotients by \(d\) mean their polynomial continuations.
\begin{table}[htbp]\centering\small
\caption{Continuous short-connector payments and their moment guards.}
\label{tab:unit-ordinary:three-low-certificates}
\begin{tabular}{@{}llc@{}}\toprule
Function & Positive factor & Degree\\\midrule
\((-F_{\ge5,3}-d/50)/d\) & \(m^4\) & 10\\
\(6/5-s_0,\ 9/8-e_0,\ e_0\) & \(m^2\) & 4 each\\
\(E_+/d,\ (E_+-E_-)/d\) & \(m^4\) & 9 each\\
\(81/256-v_0\) & \(1\) & 4\\
\(4/7-kC_0\) & \(m\) & 2\\
\(7/20-k\theta z_0\) & \(m\) & 1\\
\(1/2-k\theta^3-k^2v_0/2\) & \(m^2\) & 4\\
\((-F_{42}-d)/d\) & \(m_4^4\) & 7\\
\((-F_{43}-6d)/d\) & \(m_4^2\) & 5\\
\(7/10-s_4\) & \(m_4^2\) & 2\\\bottomrule
\end{tabular}
\end{table}
These are direct applications of \eqref{eq:bernstein-coefficients};
the exact coefficients and inverse expansions accompany the manuscript.
Also \(k\theta^2<4/7\) controls the negative part of \(a_1\),
and \(kz_0<6/7\) controls \(\delta_-\). The table therefore proves
all the moment guards used in both directions, including the endpoints.

\paragraph{Every low-domain connector of length at least four.}
For both terminal packets, \(\|P_n\|\le1\) for all \(n\ge3\).
Here is a finite-index proof. For \(b\ge5,n=3\), the real and
normalized imaginary parts are \(1-x-kx^2B\), \(1-kxA\).
They are nonnegative and bounded above by \(1-x\), \(1+kx^3\).
Hence
\[
 |P_3|^2\le1-x[1-x-2kx^3-k^2x^6]\le1;
\]
the bracket is at least \(23/256\) on \(x\le1/2,k\le3/2\).
For \(b=4\), those two components are
\(1-x-kx^3\), \(1-kdx\), between zero and the components of
\(G_3\). In both cases \(k\sqrt{A^2+xB^2}\le1\).
At \(n=4\), use
\(|P_4|\le(1-x)\sqrt{1+x}+x^2\le1\).
At \(n=5\), use
\(|P_5|\le1-x/2+(23/20)x^2\le1\).
For \(n\ge6\), the finite geometric formula gives
\(|P_n|\le(1+x^3)/\sqrt{1+x}+x^3\le1\), since
\(x(1+\sqrt{1+x})\le143/160<1\).

The same data imply the following complete loss and predecessor table:
\begin{equation}\label{eq:unit-ordinary:three-long-data}
\begin{array}{c|c|c|c|c}
 &h&\gamma_j\ (j\ge3)&\|Q_n\|\ (n\ge3)&\varepsilon/d\\\hline
 b=4&2/5&\ge hd&\le1+6d/5&[0,7/10]\\
 b\ge5&1/3&\ge hd&\le1+2d&[-1/10,9/8]
\end{array}
\end{equation}
Indeed \(\gamma_3\ge d\), \(\gamma_4>3d/4\).
For \(j\ge5\), \eqref{eq:unit-ordinary:three-losses} gives respectively
\[
 \gamma_j/d\ge\kappa-kd^2-k\theta^3\ge2/5,
 \qquad\gamma_j/d\ge\kappa-(4/7)d-k\theta^2z_0\ge1/3.
\]
For the latter bound \((D_j)_+\le\theta^2z_0d\) for every
\(j\ge4\); the four residue formulas in
\eqref{eq:unit-ordinary:three-moments} prove this also when \(A<0\).
For \(Q_n\), use its defining sum and separate \(n=3,4\) from
\(n\ge5\), so the sole \(G_2\) factor is charged \(6/5\).
The largest extra coefficients are respectively
\(231/200,1799/1600,63/64\) for \(b=4\), and
\(381/200,147/80,129/80\) for \(b\ge5\).

Retain the last two losses in
\eqref{eq:unit-ordinary:three-losses}. With \(B_*=1+\varepsilon-hd\),
the finite recursion gives
\[
 0\le T^LF/m^2\le Z=B_*q-hd+N_0,
 \quad N_0=K'Q_{L-1}(K')g',\quad \E N_0\le-hd.
\]
Weighted AM--GM at the genuine unit root gives
\[
 pT^LF\le(c_*+N/2)^2,\quad
 c_*=m\sqrt{B_*}-\frac{mhd}{2\sqrt{B_*}},\quad
 N=\frac m{\sqrt{B_*}}N_0.
\]
The square base is half a sum of nonnegative factors. The following
data concern this same complete vector, including its mean:
\begin{equation}\label{eq:unit-ordinary:three-long-payments}
\begin{array}{c|c|c|c|c|c}
 &B_*&m/\sqrt{B_*}&\|\frac m{\sqrt{B_*}}Q_{L-1}g'\|^2
 &\E N&c_*\\\hline
 b=4&[1-2d/5,1+3d/10]&\ge2/3&\le d&\le-4d/15&\le1-5d/4\\
 b\ge5&[1-13d/30,1+19d/24]&\ge5/8&\le9d/8&\le-5d/24&\le1-4d/3
\end{array}
\end{equation}
For the norm column, use respectively
\[
 m^2(1+6d/5)^2\le1-3d/5\le B_*;
 \quad m^2(1+2d)^2\le1+3d/4-(407/64)d^2
                  \le(9/8)(1-13d/30).
\]
The latter difference is bounded below by
\(1/8-(99/80)d+(407/64)d^2>0\).
The other columns follow from these \(B_*\)-intervals and
\(\sqrt{1+x}\le1+x/2\).
Substitution in \eqref{eq:unit-ordinary:three-signed} proves the
rates \(77/15\) and \(31/8\) in the final table below.

\paragraph{High-domain terminal means.}
Assume \(1/5\le d\le1/3\). The actual canonical mean is
\(s_b=\E(h_3f_b)=\E(f_2f_{b+1})\).
Let \(u_j=\E K^j\one\). The finite convolution gives
\[
 s_b=1-(2-d)\beta_{b+1}+u_{b+1}+u_{b+3}
                       -\sum_{j=3}^b\beta_js_{b-j}.
\]
Direct expansion yields
\[
\begin{gathered}
 s_0=1,\ s_1=1-d+\delta_2,\ s_2=(1-d)^2,\quad
 s_3=1-3d+d^2+\delta_2-\delta_3,\\
 s_4=1-4d+2d^2+2(1-d)\delta_2,\\
 s_5=1-5d+4d^2-d^3+(3-2d)\delta_2-\delta_3+\delta_4,\\
 s_6=1-6d+7d^2-2d^3+4(1-d)\delta_2+\delta_2^2-2(1-d)\delta_3.
\end{gathered}
\]
Consequently
\begin{equation}\label{eq:unit-ordinary:three-high-means}
 s_4\le M_4:=1-(51/16)d+(19/16)d^2,
 \quad 2M_4\le1-9d/10,\quad
 s_b\le(1-d)/2\quad(b\ge5).
\end{equation}
For the last bound at \(b=5\), use
\(s_5-(1-d)/2\le1/2-(105/32)d+(51/16)d^2-d^3
\le-147/4000\).
At \(b=6\), Cauchy--Schwarz in the spectral measure, together with
\(\delta_2\ge d^2\), gives \(\delta_3\ge d\delta_2\), and hence
\(s_6-(1-d)/2\le1/2-(31/8)d+(19/4)d^2\le-17/200\).
For every \(b\ge7\), the base in the finite convolution is at most
\(1-(2-d)\tau+\delta_4\). Retaining its five actual nonnegative
terms \(s_0,\ldots,s_4\) gives
\[
 s_b\le1-\tau\{7-11d+4d^2+(4-2d)\delta_2-\delta_3\}+\delta_4.
\]
Keep the global ratio \(w=\delta_2/d\in[d,\theta]\).
The subtracted quantity is at least
\(d\{(1-w)[7-11d+4d^2+(4-2d-\theta)wd]-\theta^2w\}\).
Its bracket decreases in \(w\), since
\(7-11d+4d^2\ge34/9\) and \((4-2d-\theta)d\le4/3\).
At \(w=\theta\) the full subtracted quantity exceeds
\((1+d)/2+7/1500\): after
retaining the cubic contribution with \(d\ge1/5\), it is enough
to check \(-1/2+(43/12)d-(53/10)d^2\ge7/1500\),
a concave quadratic with endpoint values \(7/1500,19/180\).
This proves all finite \(b\), without monotonicity of path means.

For the exchanged direction, adjunction instead gives
\[
 \widetilde s_b=\E(f_3h_b)=\E(f_4f_{b-1}),\quad
 \E f_4^2\le U_4:=1-(147/32)d+(83/16)d^2+d^3.
\]
Here the exact square is
\(1-5d+6d^2+d^3+(1-2d)\delta_2-\delta_3+\delta_4\); use
\(\delta_2\le\theta d\), \(\delta_4\le\delta_3\).
The polynomial \(U_4\) decreases on this domain, with maximum
\(U_4(1/5)=1187/4000\). The actual finite path-norm bounds of
Section~\ref{proof:Q5-C064} give, for every \(n\ge0\),
\(\|A^n\|\le1+32d/19\); they include \(n=0,1,2\) separately.
Thus
\begin{equation}\label{eq:unit-ordinary:three-high-exchanged}
 \widetilde s_b\le(1+32d/19)U_4
                 \le105643/228000<15/32\quad(b\ge5).
\end{equation}

\paragraph{Propagation, strong rates, and restoration.}
Corollary~\ref{cor:unit-square:two-step-path} and the independent
bare-path estimate \eqref{eq:Q86:0.3} give
\[
 \E(p\bar f_2)^8\le
 \begin{cases}35/8,&1/5\le d\le1/4,\\82/15,&1/4\le d\le1/3,\end{cases}
 \qquad\E(p\bar f_n)^8\le e^{253d/32}\quad(n\ge3).
\]
Since the kernel of \(T\) is at most two,
\(pT^LF\le2(\E F)p\bar f_{L-1}\).
Together with \eqref{eq:unit-ordinary:three-high-means} and
\eqref{eq:unit-ordinary:three-high-exchanged}, this yields the
following full set of rates. Entries without an expectation are
pointwise bounds on \(H=pT^LF\).
\begin{table}[htbp]\centering\small
\caption{Actual directions and complete eighth-moment payments.}
\label{tab:unit-ordinary:three-final}
\begin{tabular}{@{}lll@{}}\toprule
Actual parameters and terminal & \(0<d\le1/5\) & \(1/5\le d\le1/3\)\\\midrule
\(b=4,L=2,h_3f_4\)&\(\E H^8\le e^{-d}\)&\(H\le1-9d/10\)\\
\(b\ge5,L=2,f_3h_b\)&\(\E H^8\le e^{-4d}\)&\(H\le15/16\)\\
\(b=4,L=3,h_3f_4\)&\(\E H^8\le e^{-6d}\)&\(\E H^8\le15/16\)\\
\(b\ge5,L=3,h_3f_b\)&\(\E H^8\le e^{-d/50}\)&\(\E H^8\le3/4\)\\
\(b=4,L\ge4,h_3f_4\)&\(\E H^8\le e^{-77d/15}\)&\(\E H^8\le e^{-d/100}\)\\
\(b\ge5,L\ge4,h_3f_b\)&\(\E H^8\le e^{-31d/8}\)&\(\E H^8\le e^{-3d/32}\)\\\bottomrule
\end{tabular}
\end{table}
For \(L=3,b\ge5\), the high-domain costs are
\((4/5)^8(35/8)<3/4\) and \((3/4)^8(82/15)<3/4\).
For \(b=4\), use \(2M_4(1/5)=41/50\),
\(2M_4(1/4)=71/128<3/5\), giving respectively
\((41/50)^8(35/8)<15/16\) and \((3/5)^8(82/15)<3/4\).
For \(L\ge4,b\ge5\), use
\(8\log(1-d)+(253/32)d\le-3d/32\).
For \(b=4\), retaining three logarithm terms gives
\[
 \frac{8\log(1-9d/10)}d+253/32
 \le113/160-(81/25)d-(243/125)d^2
 \le-1951/100000<-1/100.
\]
These inequalities include both closed intervals at \(d=1/4\).
Finally \(e^{-d/100}\le1-3d/301\le1-d/500\), and the
pointwise and constant bounds are stronger still. This proves
\eqref{eq:Q89:0.3}. Positivity restores both actual equal-pair factors.
Replacing \(K\) by \(-K\) gives the full reversal, with message
\(\overline E_rqA^L(\overline E_t\,\overline F_{b,L})\).
